\begin{result}{F.3.32}{The all-genus single-line obstruction and the full CM common kernel}{Proved by an exact polynomial certificate and the actual subgroup maps}
For every $g\geq3$, every rational prime $\ell$, every nonzero
$a\in V=\Q_\ell^{2g}$, and every vector in the actual universal
degree-six image, one has
\[
 D\in U_{g,6},\qquad D(a)\in[\mathbb L_6(V),a]
             \quad\Longrightarrow\quad q_{g,6}(D)=0.       \tag{UC.1}
\]
For $X_0=E_0\setminus\{O\}$ with $E_0:y^2=x^3-x$ over
$\Q$, every nontrivial closed procyclic subgroup of the full
geometric profinite group whose arithmetic normalizer has open
image in $G_{\Q(i)}$ is contained in a conjugate of cusp inertia.
The same assertion holds after finite arithmetic base extension.
In particular, writing
\[
 R=\bigcap_\ell\ker\bigl(\pi_1(X_{0,\overline\Q})
                    \longrightarrow
                    \pi_1(X_{0,\overline\Q})^{(\ell)}\bigr),
\]
there is no nontrivial closed procyclic subgroup of $R$ with
that open-normalizer property. This resolves the full CM
common-kernel subgroup assertion in LM6 A.22.
\end{result}
\begin{proof}
We give the exact polynomial identity, its actual invariant-image
map, and the final map on full cusp inertia.

\paragraph{1. The twelve-dimensional invariant input.}
Genus three is F.3.29. For $g\geq4$, choose
$b$ with $\langle a,b\rangle=1$, write
$V=\langle a,b\rangle\perp W$, and put $s=g-1\geq3$.
Apply the actual projector $e_S$ of F.3.30, where
$S=\operatorname{Sp}(W)\times\mathbb G_m$.
It preserves $U_{g,6}\cap P_a$ and preserves $q_{g,6}$.
Thus it suffices to treat its invariant image.

The exact input construction in F.3.30 has positive
constrained-multiplicity dimensions $3,3,1,0$ and
pure-$W$ dimension five. Denote its ordered twelve
vectors by $T_1,\ldots,T_{12}$, with the seven positive
vectors first, in multiplicities $1,2,3$, followed by
the five pure vectors. The standalone certificate
\texttt{check\_all\_genus\_single\_line.py} recomputes this
basis directly from the decorated diagrams and the two
equations (LI.3). In its integer convention the positive
tensor is $m\Psi_m(A)=-\mathcal S_8(b\otimes[A,a])$;
multiplication by the nonzero integer $m$ does not change
the rational parameter space.

For the positive blocks, there are at most three $W$-pairs,
so the source invariant-tensor independence applies for
every $s\geq3$\cite[Proposition~2.1, p.~8]{AuditMoritaSakasaiSuzuki2015v2}.
For the five pure vectors, their raw five-detector determinant is
\[
 3538944\,s^5(s-2)(s-1)^4(s+1)^3(2s-1)(2s+1)^5(2s+3).
                                                               \tag{UC.2}
\]
It is nonzero for every integer $s\geq3$. The source gives
dimension five for the pure invariant special-derivation
space in all these genera\cite[Proposition~7.2, pp.~26--27]{AuditMoritaSakasaiSuzuki2015v2}.
Consequently the twelve vectors form a basis in every
required genus, including the smallest complement genus.
The independent change-of-basis check against the five
source vectors in F.3.30 gives the same specialization map.

\paragraph{2. Seven necessary scalar equations.}
Choose two paired handles $c,d$ and $e,f$ in $W$.
For $u,v,w\in V$, use the actual Johnson derivation
\[
 \delta_{uvw}(z)=\langle z,u\rangle[v,w]
              +\langle z,v\rangle[w,u]
              +\langle z,w\rangle[u,v].
\]
Let $R$ be exactly the five-term tensor contraction (SL.3).
Let $\operatorname{ES}_7$ contract the first two tensor
positions and take the cyclic class of the remaining seven
positions, with the conventions of F.3.29.
For the two six-position matchings
\[
 C_1=((1,2),(3,5),(4,6)),\qquad
 C_2=((1,2),(3,6),(4,5)),
\]
define a literal functional on the cyclic quotient by
\[
 \lambda_{C,t}([v_1\otimes\cdots\otimes v_7])
   =\sum_{j=0}^6
      \left(\prod_{(p,q)\in C}
         \langle v_{p+j},v_{q+j}\rangle\right)
         \langle v_{7+j},t\rangle,                         \tag{UC.3}
\]
where the subscripts are read modulo seven.
The complete cyclic sum makes this a well-defined
functional, including on periodic words.

In the following order, take the seven rows
\[
\begin{array}{c|c|c}
 \text{row}&\delta&\text{functional on the output}\\\hline
 1&\delta_{abc}&\langle R([D,\delta]),d\rangle\\
 2&\delta_{acd}&\langle R([D,\delta]),b\rangle\\
 3&\delta_{cde}&\langle R([D,\delta]),f\rangle\\
 4&\delta_{abc}&\lambda_{C_1,d}(\operatorname{ES}_7([D,\delta]))\\
 5&\delta_{abc}&\lambda_{C_2,d}(\operatorname{ES}_7([D,\delta]))\\
 6&\delta_{acd}&\lambda_{C_1,b}(\operatorname{ES}_7([D,\delta]))\\
 7&\delta_{cde}&\lambda_{C_1,f}(\operatorname{ES}_7([D,\delta])).
\end{array}                                                \tag{UC.4}
\]
Every row vanishes on the actual $U_{g,6}$: the actual
normalizer and strict geometric-image maps of F.3.24 give
$[U_{g,6},J_1]\subset J_7$; the ES and $R$ maps vanish
on that actual image. For $R$, the primary input is
MSS2026, Proposition~4.1, whose proof is presented there
as a sketch\cite[pp.~10--11]{AuditMoritaSakasaiSuzuki2026v1}.
Its use in any genus below the source's stable range is
justified by the exact extension map: extend the boundary
mapping class by the identity to a larger surface, extend
its derivation by zero on the added handles, and use
$R_G(i^{\otimes9}T)=iR_g(T)$. Injectivity of $i$ transfers
the asserted vanishing. The same actual stabilization and
cyclic projection map for ES was proved in F.3.24.

Let $M_7(s)$ be these seven rows on the twelve ordered
inputs, and let $K(s)$ be the five raw complete contractions
with matching indices $1,2,3,8,14$. The matrix entries are
literal polynomials in $s$. To compute them, use
\[
 T_D=\sum_i(x_i\otimes D(y_i)-y_i\otimes D(x_i)),
       \qquad D(z)=\langle\text{first tensor slot},z\rangle
                         \text{ on }T_D.
\]
Thus $T_{\delta_{uvw}}=-\operatorname{Alt}(u,v,w)$.
The two compositions in $[D,\delta]$ are formed by
substitution with this ordered contraction, with a minus
sign on $\delta\circ D$.

In a resulting fully contracted diagram, each component
is a path with two specified endpoints or a closed
alternating cycle. A path is zero unless its endpoints
have the same actual handle; a $W$-path with endpoint
$a$ or $b$ is zero. Each closed component has one
independently free handle and its two orientations give
the signed factor $2s$. Distinct closed components may
choose the same handle: their choices remain independent.
This proves the polynomial evaluation used in the
certificate, including diagrams with parallel cup and
cap edges. There are at most four $W$-edges, so the
orientation sum is finite with at most sixteen terms.
Every specified fixed endpoint in (UC.4) is among the
two displayed handles of $W$; all rows are defined for
every $s\geq3$.

\paragraph{3. A literal all-parameter row identity.}
For the two actual geometric brackets of F.3.29, put
\[
 J=\begin{pmatrix}
 -104&-288\\-26&-72\\104&288\\32&192\\-32&-192
 \end{pmatrix},\qquad
 L=\begin{pmatrix}-1&4&0&0&0\\1&0&1&0&0\\0&0&0&1&1\end{pmatrix}.
\]
Their ranks are two and three and $LJ=0$, so
$\ker L=\operatorname{im}J$.
Set
\[
                 d_0(s)=9(2s+3)(4s^2+6s-1).
\]
The following explicit $3$ by $7$ polynomial matrix
$N_0$ satisfies
\[
                    N_0(s)M_7(s)=d_0(s)LK(s).              \tag{UC.5}
\]
To specify every entry without a wide matrix, its rows are:
\begin{align*}
 (N_0)_{11}&=-2s(2s+3)(128s^4+528s^3+776s^2+465s+101),\\
 (N_0)_{12}&=-4s(2s+3)(16s^4-48s^3-305s^2-402s-152),\\
 (N_0)_{13}&=2s(s-1)(2s+3)(128s^3+464s^2+550s+199),\\
 (N_0)_{14}&=-3s(576s^4+2720s^3+4684s^2+3546s+1029),\\
 (N_0)_{15}&=-6s(s+2)(4s+3)(4s+5)(8s+7),\\
 (N_0)_{16}&=3s(2s+1)(2s+3)(80s^2+240s+187),\\
 (N_0)_{17}&=24s(s-1)(2s+3)(4s^2+6s-1);\\[2pt]
 (N_0)_{21}&=2s(2s+1)^2(2s+3)(8s^2+25s+20),\\
 (N_0)_{22}&=2s(2s+1)(2s+3)(4s^3-14s^2-70s-61),\\
 (N_0)_{23}&=-4s(s-1)(2s+1)(2s+3)(8s^2+25s+20),\\
 (N_0)_{24}&=3s(2s+3)(72s^3+232s^2+224s+69),\\
 (N_0)_{25}&=6s(s+2)(2s+1)(2s+3)(8s+7),\\
 (N_0)_{26}&=-3s(2s+3)(40s^3+140s^2+146s+37),\\
 (N_0)_{27}&=-6s(s-1)(2s+3)(4s^2+6s-1);\\[2pt]
 (N_0)_{31}&=-6s(2s+3)(16s^3+66s^2+80s+27),\\
 (N_0)_{32}&=-6s(2s+3)(4s^3-12s^2-82s-81),\\
 (N_0)_{33}&=6s(s-1)(2s+3)(16s^2+60s+53),\\
 (N_0)_{34}&=-18s(s+2)(36s^2+102s+67),\\
 (N_0)_{35}&=-36s(s+2)^2(8s+7),\\
 (N_0)_{36}&=18s(s+2)(2s+3)(10s+13),\\
 (N_0)_{37}&=0.
\end{align*}
The standalone program recomputes all twelve inputs,
all sixty entries of $K$, all eighty-four entries of
$M_7$, and every entry of the difference in (UC.5).
All thirty-six polynomial differences are zero in
$\Q[s]$. Modular arithmetic is used only to choose
independent input rows: the input kernels are checked
against every original integer equation, and all
polynomial identities use exact integer or rational
arithmetic. Explicit integer-size bounds precede the
two array matrix products. No stored matrix, input
pickle, interpolation, or numerical root test is required.
Independent word realizations with two $W$-handles
checked all sixty detector, thirty-six $R$, and
forty-eight ES entries against the same polynomials;
these are checks of the literal maps, separate from
the polynomial identity itself.

Writing $s=3+t$ gives
\[
                   d_0(3+t)=9(9+2t)(53+30t+4t^2)>0
                           \qquad(t\geq0).
\]
There are therefore no exceptional integers $s\geq3$.
Equation (UC.5) proves, over every $\Q_\ell$,
\[
                       \ker M_7(s)\subset\ker LK(s).
                                                               \tag{UC.6}
\]
Only these seven necessary rows are used; no assertion
that they exhaust the mixed equations is needed.

\paragraph{4. Return to the actual universal quotient.}
Take $D\in(U_{g,6}\cap P_a)^S$ and write its invariant
tensor in the twelve-vector basis. Its coordinate
vector $v$ satisfies $M_7(s)v=0$, hence $LK(s)v=0$.
The raw-to-source five-coordinate map is precisely
$\operatorname{diag}(1,-1,1,-1,-1)$, so the injectivity
and span assertions of the source detector are
transported by that fixed invertible map.
The exact inclusion (LI.6) now gives $q_{g,6}(D)=0$:
after the full symplectic invariant projector, its
five coordinates are those of a linear combination
of the two actual geometric brackets, and the
source detector is injective on that invariant
space. The quotient kills these actual geometric
brackets. Finally $q_{g,6}e_S=q_{g,6}$ proves (UC.1)
for the original vector. This uses a projector on
the actual universal image; no averaged vector is
declared to be an individual Galois element.

\paragraph{5. Full profinite cuspidality and the common kernel.}
Apply the unconditional witness map of F.3.31 to a
nonperipheral closed procyclic subgroup of the full
geometric group of $X_0$. It first uses the actual
proper-homology detector of Nakamura1994, Lemma~2.1.4,
before taking any pro-prime quotient\cite[pp.~96--97]{AuditNakamura1994}.
It preserves a nonzero proper class through the
finite-cover power, arithmetic normal core, genus
raising, rational tangential descent, and filling
to a once-punctured curve. The actual nonzero
degree-six arithmetic class for that same curve is
provided by Nakamura1996, printed pp.~115--117,
including the finite-base invariance and fibre
surjection\cite{AuditNakamura1996}. The normalizer
lift and full-logarithm correction give a vector
in $U_{g,6}\cap P_a$ with nonzero quotient. Genus
three contradicts F.3.29 and every higher genus
contradicts (UC.1). Equivalently, the matrix witness
in (FC.7) contradicts (UC.6), since all seven rows
are among the necessary equations of that witness.
Thus a topological generator of the subgroup lies
in a cusp inertia subgroup, and its closed procyclic
closure lies there as well. The argument only asks
that the arithmetic normalizer contain an open
Galois subgroup; finite arithmetic base extensions
preserve this hypothesis and the actual source
quotient used in the proof.

For completeness, the final exclusion from $R$ is
an explicit map, not an inference from primewise
recognition. Write the geometric free profinite
group as $\widehat{\langle x,y\rangle}$ and choose
the inertia generator $w=[x,y]$, up to inversion.
Its maximal pro-$\ell$ image $w_\ell$ has degree-two
Magnus class $[X,Y]$ in the free rank-two pro-$\ell$
group. This degree-two module is free over $\Z_\ell$.
Equivalently, the substitutions
\[
 x\longmapsto 1+E_{12},\qquad
 y\longmapsto 1+E_{23}
       \quad\text{in }\operatorname{UT}_3(\Z/\ell^n\Z)
\]
send $w$ to $1+E_{13}$, which has exact order
$\ell^n$ for every $n$. Hence cusp inertia is
$I\simeq\widehat\Z$, and for
$\alpha=(\alpha_\ell)_\ell\in\widehat\Z$ its actual
image under the maximal pro-$\ell$ quotient is
$w_\ell^{\alpha_\ell}$. Its degree-two class is
$\alpha_\ell[X,Y]$, which is nonzero whenever
$\alpha_\ell\ne0$. Every nonzero $\alpha$ has such
a component, so
\[
                             I\cap R=1.
\]
The same equality holds for every conjugate of $I$,
because the quotient maps preserve conjugation
and $R$ is normal. A subgroup of $R$ with the
open-normalizer property is cuspidal by the
preceding paragraph and therefore trivial.
The general affine-curve extension is a separate
assertion; only the stated original CM curve and
its finite arithmetic base changes are concluded here.
\end{proof}
